% Event-timing-focused revised introduction for Overleaf.
% Citations are written as author-year placeholders for later integration.

\documentclass[11pt]{article}
\usepackage[margin=1in]{geometry}
\usepackage{amsmath}
\usepackage[version=4]{mhchem}
\usepackage[T1]{fontenc}
\usepackage{lmodern}

\title{Precession-Phase Organization of DO-Like Asian Monsoon Transitions}
\author{}
\date{}

\begin{document}
\maketitle

\section*{Introduction}

Millennial-scale climate variability (MCV) is a major expression of abrupt climate change in glacial climate systems. Its canonical form is the Dansgaard--Oeschger (DO) sequence in Greenland ice cores, where rapid stadial--interstadial transitions record large reorganizations of the North Atlantic climate system (Rasmussen et al., 2014). However, MCV is not solely a Greenland phenomenon. Marine, ice-core, and speleothem archives indicate that abrupt millennial variability occurred across multiple Pleistocene glacial cycles and was expressed in both high-latitude and low-latitude climate systems (Barker et al., 2011; Sun et al., 2021; Hodell et al., 2023). Speleothem records are especially important because their U--Th chronologies allow abrupt hydroclimate changes to be compared directly with ice-core event sequences. During the last glacial period, independently dated speleothem records show that abrupt Greenland warmings were accompanied by near-synchronous hydroclimate shifts in monsoon and mid-latitude regions (Corrick et al., 2020). Chinese speleothem weak and strong monsoon transitions can therefore be treated as precisely dated DO-like Asian monsoon expressions of MCV, while remaining distinct from Greenland-defined DO events.

A growing body of work suggests that MCV is not independent of orbital-scale boundary conditions. Sun et al. (2021) showed that MCV persisted through the Pleistocene and that its amplitude was modulated by precession, obliquity, and glacial background state. North Atlantic records likewise indicate coupling between orbital and millennial variability, with MCV amplitude varying with ice volume and orbital configuration (Hodell et al., 2023). In the monsoon domain, Thirumalai et al. (2020) used the Cheng et al. (2016) Chinese speleothem composite, together with atmospheric methane, to examine orbital modulation of millennial-scale variability. Their results placed Chinese speleothem millennial variability within the broader MCV framework and showed that orbital forcing can modulate the magnitude of millennial-scale signals, although the expression of this modulation differs among proxies. These studies establish that orbital configuration and glacial background state can shape the amplitude or envelope of MCV.

Amplitude modulation, however, is not equivalent to event pacing. A stronger millennial-scale signal under a given orbital configuration does not imply that individual abrupt transitions occur preferentially at a particular orbital phase. Previous proxy studies have clarified the synchrony of speleothem events with Greenland abrupt changes (Corrick et al., 2020), the persistence and orbital modulation of MCV amplitude through the Pleistocene (Sun et al., 2021; Hodell et al., 2023), and the modulation of millennial-scale monsoon variability in the Cheng et al. (2016) speleothem record (Thirumalai et al., 2020). What remains unresolved is an event-level timing question: are objectively identified weak and strong Asian monsoon starts randomly distributed through orbital phase, or do they cluster within preferred precession or obliquity windows? Climate-model experiments provide a dynamical motivation for this question, showing that astronomical configuration can directly influence simulated DO-like oscillations and abrupt transitions through changes in tropical hydroclimate, sea ice, and ocean circulation (Zhang et al., 2021). Yet this model-based possibility requires a proxy-based event-timing test.

Here we test whether speleothem-recorded DO-like Asian monsoon transitions are organized by orbital phase. We use the 0--640 ka sequence of weak and strong Asian monsoon starts identified by Rousseau et al. (2023) from the Cheng et al. (2016) Chinese speleothem composite. These events are analyzed as monsoon transition proxies for MCV, not as Greenland-defined DO events. We first evaluate whether the transition sequences are consistent with stationary random occurrence. We then use circular Rayleigh tests to assess phase clustering, and binned Poisson hazard generalized linear models to test whether orbital phase adds predictive information beyond LR04 ice volume and atmospheric \(\ce{CO2}\). This framework shifts the focus from the magnitude of millennial-scale variability to the timing of objectively identified monsoon transitions.

The results show that weak and strong monsoon starts are inconsistent with simple stationary randomness and exhibit significant, opposite preferences in precession phase. Poisson hazard models further indicate that precession phase remains informative after accounting for LR04 and \(\ce{CO2}\), whereas obliquity phase and other candidate forcings do not provide comparable independent contributions. These findings suggest that DO-like Asian monsoon transitions were not merely stronger or weaker under particular orbital configurations; their occurrence probability was itself organized by precession phase. We therefore interpret precession as a phase-dependent boundary condition that modulates transition likelihood together with slower ice-volume and greenhouse-gas background states.

\end{document}
